%%% Hlavní soubor. Zde se definují základní parametry a odkazuje se na ostatní části. %%%

%% Verze pro jednostranný tisk:
% Okraje: levý 40mm, pravý 25mm, horní a dolní 25mm
% (ale pozor, LaTeX si sám přidává 1in)
\documentclass[12pt,a4paper]{report}
\setlength\textwidth{145mm}
\setlength\textheight{247mm}
\setlength\oddsidemargin{15mm}
\setlength\evensidemargin{15mm}
\setlength\topmargin{0mm}
\setlength\headsep{0mm}
\setlength\headheight{0mm}
% \openright zařídí, aby následující text začínal na pravé straně knihy
\let\openright=\clearpage

%% Pokud tiskneme oboustranně:
% \documentclass[12pt,a4paper,twoside,openright]{report}
% \setlength\textwidth{145mm}
% \setlength\textheight{247mm}
% \setlength\oddsidemargin{14.2mm}
% \setlength\evensidemargin{0mm}
% \setlength\topmargin{0mm}
% \setlength\headsep{0mm}
% \setlength\headheight{0mm}
% \let\openright=\cleardoublepage

%% Vytváříme PDF/A-2u
\usepackage[a-2u]{pdfx}
\usepackage{listings}
\def\inline{\lstinline[basicstyle=\ttfamily,keywordstyle={}]}

%% Přepneme na českou sazbu a fonty Latin Modern
\usepackage[czech]{babel}
\usepackage{lmodern}
\usepackage[T1]{fontenc}
\usepackage{textcomp}

%% Použité kódování znaků: obvykle latin2, cp1250 nebo utf8:
\usepackage[utf8]{inputenc}

%%% Další užitečné balíčky (jsou součástí běžných distribucí LaTeXu)
\usepackage{amsmath}        % rozšíření pro sazbu matematiky
\usepackage{amsfonts}       % matematické fonty
\usepackage{amsthm}         % sazba vět, definic apod.
\usepackage{bbding}         % balíček s nejrůznějšími symboly
			    % (čtverečky, hvězdičky, tužtičky, nůžtičky, ...)
\usepackage{bm}             % tučné symboly (příkaz \bm)
\usepackage{graphicx}       % vkládání obrázků
\usepackage{fancyvrb}       % vylepšené prostředí pro strojové písmo
\usepackage{indentfirst}    % zavede odsazení 1. odstavce kapitoly
\usepackage{natbib}         % zajištuje možnost odkazovat na literaturu
			    % stylem AUTOR (ROK), resp. AUTOR [ČÍSLO]
\usepackage[nottoc]{tocbibind} % zajistí přidání seznamu literatury,
                            % obrázků a tabulek do obsahu
\usepackage{icomma}         % inteligetní čárka v matematickém módu
\usepackage{dcolumn}        % lepší zarovnání sloupců v tabulkách
\usepackage{booktabs}       % lepší vodorovné linky v tabulkách
\usepackage{paralist}       % lepší enumerate a itemize
\usepackage[usenames]{xcolor}  % barevná sazba

%%% Údaje o práci

\def\NazevSkoly{Gymnázium, Praha 6, Arabská 14}
% Název oboru včetně počátečního 'Obor'.
\def\NazevOboru{Obor programování}

% Název práce v jazyce práce (přesně podle zadání)
\def\NazevPrace{<NazevPrace>}

% Název práce v angličtině
\def\NazevPraceEN{<NazevPraceEN>}

% Jména autorů
% Abecedně podle příjmení
\def\AutorPrace{Tomáš~Černý}

% Rok odevzdání
\def\RokOdevzdani{2020}
% Měsíc odevzdání
\def\MesicOdevzdani{<MesicOdevzdani>}

% Vedoucí práce: Jméno a příjmení s~tituly
\def\Vedouci{Mgr. Jan Lána}

% Nepovinné poděkování (vedoucímu práce, konzultantovi, tomu, kdo
% zapůjčil software, literaturu apod.)
\def\Podekovani{%
}

% Abstrakt (doporučený rozsah cca 80-200 slov; nejedná se o zadání práce)
\def\Abstrakt{%
<Abstrakt>
}
\def\AbstraktEN{%
<AbstraktEN>
}

% 3 až 5 klíčových slov (doporučeno), každé uzavřeno ve složených závorkách
\def\KlicovaSlova{%
{klíčová} {slova}
}
\def\KlicovaSlovaEN{%
{key} {words}
}

%% Balíček hyperref, kterým jdou vyrábět klikací odkazy v PDF,
%% ale hlavně ho používáme k uložení metadat do PDF (včetně obsahu).
%% Většinu nastavítek přednastaví balíček pdfx.
\hypersetup{unicode}
\hypersetup{breaklinks=true}

%% Definice různých užitečných maker (viz popis uvnitř souboru)
%%% Tento soubor obsahuje definice různých užitečných maker a prostředí %%%
%%% Další makra připisujte sem, ať nepřekáží v ostatních souborech.     %%%

%%% Drobné úpravy stylu

% Tato makra přesvědčují mírně ošklivým trikem LaTeX, aby hlavičky kapitol
% sázel příčetněji a nevynechával nad nimi spoustu místa. Směle ignorujte.
\makeatletter
\def\@makechapterhead#1{
  {\parindent \z@ \raggedright \normalfont
   \Huge\bfseries \thechapter. #1
   \par\nobreak
   \vskip 20\p@
}}
\def\@makeschapterhead#1{
  {\parindent \z@ \raggedright \normalfont
   \Huge\bfseries #1
   \par\nobreak
   \vskip 20\p@
}}
\makeatother

% Toto makro definuje kapitolu, která není očíslovaná, ale je uvedena v obsahu.
\def\chapwithtoc#1{
\chapter*{#1}
\addcontentsline{toc}{chapter}{#1}
}

% Trochu volnější nastavení dělení slov, než je default.
\lefthyphenmin=2
\righthyphenmin=2

% Zapne černé "slimáky" na koncích řádků, které přetekly, abychom si
% jich lépe všimli.
% \overfullrule=1mm

%%% Makra pro definice, věty, tvrzení, příklady, ... (vyžaduje baliček amsthm)

\theoremstyle{plain}
\newtheorem{veta}{Věta}
\newtheorem{lemma}[veta]{Lemma}
\newtheorem{tvrz}[veta]{Tvrzení}

\theoremstyle{plain}
\newtheorem{definice}{Definice}

\theoremstyle{remark}
\newtheorem*{dusl}{Důsledek}
\newtheorem*{pozn}{Poznámka}
\newtheorem*{prikl}{Příklad}

%%% Prostředí pro důkazy

\newenvironment{dukaz}{
  \par\medskip\noindent
  \textit{Důkaz}.
}{
\newline
\rightline{$\square$}  % nebo \SquareCastShadowBottomRight z balíčku bbding
}

%%% Prostředí pro sazbu kódu, případně vstupu/výstupu počítačových
%%% programů. (Vyžaduje balíček fancyvrb -- fancy verbatim.)

\DefineVerbatimEnvironment{code}{Verbatim}{fontsize=\small, frame=single}

%%% Prostor reálných, resp. přirozených čísel
\newcommand{\R}{\mathbb{R}}
\newcommand{\N}{\mathbb{N}}

%%% Užitečné operátory pro statistiku a pravděpodobnost
\DeclareMathOperator{\pr}{\textsf{P}}
\DeclareMathOperator{\E}{\textsf{E}\,}
\DeclareMathOperator{\var}{\textrm{var}}
\DeclareMathOperator{\sd}{\textrm{sd}}

%%% Příkaz pro transpozici vektoru/matice
\newcommand{\T}[1]{#1^\top}

%%% Vychytávky pro matematiku
\newcommand{\goto}{\rightarrow}
\newcommand{\gotop}{\stackrel{P}{\longrightarrow}}
\newcommand{\maon}[1]{o(n^{#1})}
\newcommand{\abs}[1]{\left|{#1}\right|}
\newcommand{\dint}{\int_0^\tau\!\!\int_0^\tau}
\newcommand{\isqr}[1]{\frac{1}{\sqrt{#1}}}

%%% Vychytávky pro tabulky
\newcommand{\pulrad}[1]{\raisebox{1.5ex}[0pt]{#1}}
\newcommand{\mc}[1]{\multicolumn{1}{c}{#1}}


%% Titulní strana a různé povinné informační strany
\begin{document}
%%% Titulní strana práce a další povinné informační strany

%%% Titulní strana práce

\pagestyle{empty}
\hypersetup{pageanchor=false}

\begin{center}

{\LARGE\bfseries\NazevSkoly}

\vspace{-22mm}
\vfill

{\LARGE\NazevOboru}

\vfill

\centerline{\mbox{\includegraphics[width=100mm]{../img/logo-ga-new.pdf}}}

\vspace{-8mm}
\vfill

{\bf\Large ROČNÍKOVÝ PROJEKT}

\vfill

{\Large \AutorPrace}

\vspace{15mm}

{\LARGE\bfseries\NazevPrace}

% \vfill
% \Katedra
% \vfill

% \begin{tabular}{rl}

% Vedoucí bakalářské práce: & \Vedouci \\
% \noalign{\vspace{2mm}}

\vfill

\MesicOdevzdani \ \RokOdevzdani

\end{center}

\newpage

%%% Následuje vevázaný list -- kopie podepsaného "Zadání bakalářské práce".
%%% Toto zadání NENÍ součástí elektronické verze práce, nescanovat.

%%% Strana s čestným prohlášením k bakalářské práci

\openright
\hypersetup{pageanchor=true}
\pagestyle{plain}
\pagenumbering{roman}
\vglue 0pt plus 1fill

\noindent Prohlašuji, že jsem jediným autorem tohoto projektu, všechny citace
jsou řádně označené a všechna použitá literatura a další zdroje jsou v práci
uvedené.  Tímto dle zákona 121/2000 Sb. (tzv. Autorský zákon) ve znění
pozdějších předpisů uděluji bezúplatně škole Gymnázium, Praha 6, Arabská14
oprávnění k výkonu práva na rozmnožování díla (§ 13) a práva na sdělování díla
veřejnosti (§ 18) na dobu časově neomezenou a bez omezení územního rozsahu.

\vspace{10mm}

\hbox{\hbox to 0.5\hsize{%
V ........ dne ............
\hss}\hbox to 0.5\hsize{%
Tomáš Černý ............
\hss}}

\vspace{20mm}

% \newpage
%
% %%% Poděkování
%
% \openright
%
% \noindent
% \Podekovani

\newpage

%%% Povinná informační strana bakalářské práce

\openright

\vbox to 0.5\vsize{
\setlength\parindent{0mm}
\setlength\parskip{5mm}

Název práce:
\NazevPrace

Autoři:
\AutorPrace

% TODO: Je Lána vedoucí?
% Vedoucí práce:
% \Vedouci, \KatedraVedouciho

Abstrakt:
\Abstrakt

% Klíčová slova:
% \KlicovaSlova

\vss}\nobreak\vbox to 0.49\vsize{
\setlength\parindent{0mm}
\setlength\parskip{5mm}

% Opakování v angličtině.

Title:
\NazevPraceEN

Authors:
\AutorPrace

Abstract:
\AbstraktEN

\vss}

\newpage

\openright
\pagestyle{plain}
\pagenumbering{arabic}
\setcounter{page}{1}


%%% Strana s automaticky generovaným obsahem bakalářské práce

\tableofcontents

%%% Jednotlivé kapitoly práce jsou pro přehlednost uloženy v samostatných souborech
\chapter{Úvod} \label{uvod}

\noindent
Tento dokument se zabývá koncepcí a implementací informačního systému pro uzavřená
parkoviště s automatickým rozpoznáváním SPZ bez použití drahého kamerového hardware.
Inspirací pro projekt bylo, že se po autorovi při výjezdu z parkoviště nechtěl
parkovací lístek, nýbrž vozidlo bylo puštěno na základě SPZ.

Projekt se nezabývá operací se závorou a platebním terminálem, neb by to zvyšovalo
obtížnost už tak obtížného projektu a vyžadovalo by to poměrně vysoké
náklady na hardware.

Projekt zároveň slouží k vyzkoušení si moderních webových technologií
a frameworků s přesahem k mobilnímu vývoji a počítačovému vidění.

\section*{Zadání} \label{zadani}

\noindent
Parkovací systém má mít následující funkcionalitu:

\begin{enumerate}
  \setlength\itemsep{0.05em}
  \item Naskenování SPZ.
  \item Vytváření dostatečně flexibilních pravidel pro většinu použití.
  \begin{enumerate}
    \setlength\itemsep{0.05em}
      \item Od času A do B za tarif X/jednotka času.
      \item Různý provoz o svátcích a víkendech.
      \item Limity na hodiny zdarma.
      \item Filtrování vozidel -- pro různá vozidla mohou platit jiná pravidla.
    \end{enumerate}
  \item Statistiky počtu aut a výdělku s grafy.
\end{enumerate}

%%% Fiktivní kapitola s ukázkami sazby

\chapter{Kapitola} \label{kapitola1}

<Content>

\citep[viz][]{examplecit}


\chapter*{Závěr}
\addcontentsline{toc}{chapter}{Závěr}

Výsledný projekt splňuje celé zádání kromě jednoho bodu -- konkrétně se jedná o
různý provoz
o svátích a podobných dnech (viz \ref{missing1}). Implementace takovéto funkcionality
vyžaduje implementovat kalendář.
Co se týče rozpoznávání SPZ, statistik, samotných pravidel a filtrů vozidel,
tak zde je vše implementováno a funguje skvěle.
Při prohlížení záznamů parkování lze nahlédnout na výřez SPZ z pořízeného snímku.
Navíc lze pravidla a filtry simulovat pro libovolné vozidlo.

Vývoj byl díky vhodně zvoleným technologiím poměrně rychlý a bezbolestný.
V jeho průběhu nedošlo k žádnému backtrackování kvůli předchozím rozhodnutím.
Projekt je bez velkých obtíží rozšiřitelný a pozměnitelný.
Dalším rozšířením, aby projekt byl plnohodnotný parkovací systém, by byla
integrace s platebním terminálem a závorou.


%%% Seznam použité literatury
%%% Seznam použité literatury (bibliografie)
%%%
%%% Pro vytváření bibliografie používáme bibTeX. Ten zpracovává
%%% citace v textu (např. makro \cite{...}) a vyhledává k nim literaturu
%%% v souboru literatura.bib.
%%%
%%% Příkaz \bibliographystyle určuje, jakým stylem budou citovány odkazy
%%% v textu. V závorce je název zvoleného souboru .bst. Styly plainnat
%%% a unsrt jsou standardní součástí latexových distribucí. Styl czplainnat
%%% je dodáván s touto šablonou a bibTeX ho hledá v aktuálním adresáři.

\bibliographystyle{czplainnat}    %% Autor (rok) s českými spojkami
% \bibliographystyle{plainnat}    %% Autor (rok) s anglickými spojkami
% \bibliographystyle{unsrt}       %% [číslo]

\renewcommand{\bibname}{Seznam použité literatury}

%%% Vytvoření seznamu literatury. Pozor, pokud jste necitovali ani jednu
%%% položku, seznam se automaticky vynechá.

\bibliography{literatura}

%%% Kdybyste chtěli bibliografii vytvářet ručně (bez bibTeXu), lze to udělat
%%% následovně. V takovém případě se řiďte normou ISO 690 a zvyklostmi v oboru.

% \begin{thebibliography}{99}
%
% \bibitem{lamport94}
%   {\sc Lamport,} Leslie.
%   \emph{\LaTeX: A Document Preparation System}.
%   2. vydání.
%   Massachusetts: Addison Wesley, 1994.
%   ISBN 0-201-52983-1.
%
% \end{thebibliography}


%%% Obrázky v bakalářské práci
%%% (pokud jich je malé množství, obvykle není třeba seznam uvádět)
\listoffigures

%%% Tabulky v bakalářské práci (opět nemusí být nutné uvádět)
%%% U matematických prací může být lepší přemístit seznam tabulek na začátek práce.
\listoftables

%%% Použité zkratky v bakalářské práci (opět nemusí být nutné uvádět)
%%% U matematických prací může být lepší přemístit seznam zkratek na začátek práce.
% \chapwithtoc{Seznam použitých zkratek}

%%% Přílohy k bakalářské práci, existují-li. Každá příloha musí být alespoň jednou
%%% odkazována z vlastního textu práce. Přílohy se číslují.
%%%
%%% Do tištěné verze se spíše hodí přílohy, které lze číst a prohlížet (dodatečné
%%% tabulky a grafy, různé textové doplňky, ukázky výstupů z počítačových programů,
%%% apod.). Do elektronické verze se hodí přílohy, které budou spíše používány
%%% v elektronické podobě než čteny (zdrojové kódy programů, datové soubory,
%%% interaktivní grafy apod.). Elektronické přílohy se nahrávají do SISu a lze
%%% je také do práce vložit na CD/DVD. Povolené formáty souborů specifikuje
%%% opatření rektora č. 72/2017.
% \appendix
% \chapter{Přílohy}

% \openright
\end{document}
